\section{Introduction}
\label{sec:intro}

\begin{guidance}[Introduction (10 marks), about half a page]
The marker wants: the context and goal of the project; your \textbf{selected
game} and \textbf{GUI library}, with the reason for each choice; your main
\textbf{design choices}; the \textbf{task allocation} between the two of you; and
the \textbf{structure} of the report. Keep it to about half a page (brief,
p.~6). Open with the problem, not with the technology. Expand every acronym at
first use, e.g.\ ``finite-state machine (FSM)''.
\end{guidance}

In networked and serial gaming architectures, placing game logic or rule enforcement on the client exposes the system to desynchronization, packet manipulation, and cheating. If the client is trusted to report hits, misses, or board updates, state integrity cannot be guaranteed. As specified in the EEE3095S project brief~\cite{brief}, the solution is an authoritative referee model: a dedicated microcontroller owns the game state, validates every client transaction against strict game rules, and returns unambiguous evaluations across an asynchronous serial link. A fundamental property enforced throughout this design is \textbf{strict state invariance}: any request that is malformed, out-of-order, or invalid according to the current game phase is rejected immediately via an explicit negative acknowledgement (NACK), guaranteeing that internal server state remains completely uncorrupted.

To rigorously exercise this client-untrusted architecture, \textbf{Battleship} was selected over baseline games such as Tic-Tac-Toe. While Tic-Tac-Toe features full public board observability, Battleship inherently requires \textit{hidden information} (fog of war). This establishes a demanding test of server authority and secrecy: the microcontroller must maintain both players' secret fleet layouts in internal SRAM, evaluate incoming coordinates, and emit only localized outcomes (\texttt{HIT}, \texttt{MISS}, or \texttt{SUNK}) without ever leaking unhit ship coordinates until match conclusion. The PC client is built using \textbf{PyQt6}~\cite{guilib}, chosen for its robust signal-and-slot architecture, modern styling, and native multi-threading support via \texttt{QThread}. This enables a strictly thin client: all rule evaluations are prohibited on the PC, user inputs are locked pending microcontroller confirmation, and serial communication runs on an isolated background thread via pySerial~\cite{pyserial} to prevent graphical freezes during serial transfer.

Key design choices shape the system into a resilient embedded application. The firmware targeting the STM32F051 ARM Cortex-M0 is strictly modularized into protocol framing, transactional game logic, asynchronous hardware control, and an autonomous probability density heatmap artificial intelligence (AI) opponent. Framing employs newline-delimited ASCII strings protected by an inter-byte silence timeout and memory-safe boundary checks. Peripheral integration features a debounced hardware reset push-button, dual-colour turn indicators, and a dynamic five-LED fleet health bar. Finally, system correctness is verified using an automated, unattended test harness built on \texttt{pytest}~\cite{pytest} that drives the physical hardware over serial with automated port discovery.

\textbf{Project execution and task allocation.} This project was completed individually by Scott Faurholm (FRHSCO001). Consequently, all engineering deliverables---including the protocol specification, STM32F051 embedded firmware (FSM, atomic placement validation, hardware peripherals, and probability heatmap AI), the PyQt6 thin GUI client, the automated \texttt{pytest} test harness, and this report---were designed, implemented, and validated independently by the author.

\textbf{Structure.} The remainder of this report is organized as follows: Section~\ref{sec:protocol} defines the communication protocol, framing rules, command set, and departures from the reference specification. Section~\ref{sec:design} presents the system architecture, embedded FSM, and circuit schematic. Section~\ref{sec:impl} analyzes key implementation excerpts including atomic placement and the heatmap AI. Section~\ref{sec:authority} demonstrates server authority and proof of state invariance. Section~\ref{sec:robust} evaluates framing resilience, error responses, and garbage recovery. Section~\ref{sec:valid} presents automated test results and hardware validation. Finally, Section~\ref{sec:conclusion} concludes with limitations and proposed enhancements.
